\documentclass{beamer}
\usepackage{tikz,amsmath,hyperref,graphicx,stackrel,animate,amssymb}
\usetikzlibrary{positioning,shadows,arrows,shapes,calc}
\newcommand{\argmax}{\operatornamewithlimits{argmax}}
\newcommand{\argmin}{\operatornamewithlimits{argmin}}
\mode<presentation>{\usetheme{Frankfurt}}
\AtBeginSection[]
{
  \begin{frame}<beamer>
    \frametitle{Outline}
    \tableofcontents[currentsection,currentsubsection]
  \end{frame}
}
\title{Lecture 13: Temporal and Spatial Aliasing}
\author{Mark Hasegawa-Johnson}
\date{ECE 401: Signal and Image Analysis}
\begin{document}

% Title
\begin{frame}
  \maketitle
\end{frame}

% Title
\begin{frame}
  \tableofcontents
\end{frame}

%%%%%%%%%%%%%%%%%%%%%%%%%%%%%%%%%%%%%%%%%%%%
\section[Review]{Review: DFT}
\setcounter{subsection}{1}

\begin{frame}
  \frametitle{Review: DFT}

  The DFT (discrete Fourier transform) of any signal is
  $X[k]$, given by
  \begin{align*}
    X[k] &= \sum_{n=0}^{N-1} x[n]e^{-j\frac{2\pi kn}{N}}\\
    x[n] &= \frac{1}{N}\sum_0^{N-1} X[k]e^{j\frac{2\pi kn}{N}}
  \end{align*}
  Particular useful examples include:
  \begin{align*}
    f[n]=\delta[n] &\leftrightarrow F[k]=1\\
    g[n]=\delta\left[(\!(n-n_0)\!)_N\right] &\leftrightarrow G[k]=e^{-j\frac{2\pi kn_0}{N}}
  \end{align*}
\end{frame}

\begin{frame}
  \frametitle{Properties of the DTFT}

  Properties worth knowing  include:
  \begin{enumerate}
    \setcounter{enumi}{-1}
  \item Periodicity: $X[k+N]=X[k]$
  \item Conjugate symmetric: $x[n]$ Real $\leftrightarrow X[k]=X^*[-k]$
  \item Linearity:
    \[z[n]=ax[n]+by[n]\leftrightarrow Z[k]=aX[k]+bY[k]
    \]
  \item Time Shift: $x[n-n_0]\leftrightarrow e^{-j\frac{2\pi
      kn_0}{N}}X[k]$, paying attention to the fact that $x[n]$ is
    periodic in time
  \item Frequency Shift: $e^{j\frac{2\pi k_0 n}{N}}x[n]\leftrightarrow
    X[k-k_0]$, paying attention to the fact that $X[k]$ is periodic
    in frequency
  \end{enumerate}
\end{frame}

%%%%%%%%%%%%%%%%%%%%%%%%%%%%%%%%%%%%%%%%%%%%
\section[Periodic in Time]{Sampled in Frequency $\leftrightarrow$ Periodic in Time}
\setcounter{subsection}{1}

\begin{frame}
  \frametitle{Try the quiz!}

  Today's quiz is another circular convolution problem, because that
  might be the best way to think about periodicity in time.  Go to the
  course web page and try it!
\end{frame}

\begin{frame}
  \frametitle{$x[n]$ periodic, $X[k]=NX_k$}

  If $x[n]=x[n+N]$, then its Fourier series is
  \begin{align*}
    X_k &= \frac{1}{N}\sum_{n=1}^{N-1} x[n]e^{-j\frac{2\pi kn}{N}}\\
    x[n] &= \sum_{k=0}^{N-1} X_k e^{j\frac{2\pi kn}{N}},
  \end{align*}
  and its DFT is
  \begin{align*}
    X[k] &= \sum_{n=1}^{N-1} x[n]e^{-j\frac{2\pi kn}{N}}\\
    x[n] &= \frac{1}{N}\sum_{k=0}^{N-1} X[k] e^{j\frac{2\pi kn}{N}}
  \end{align*}
\end{frame}
  
\begin{frame}
  \frametitle{Delayed impulse wraps around}

  \begin{displaymath}
    \delta\left[(\!(n-n_0)\!)_N\right] \leftrightarrow e^{-j\frac{2\pi kn_0}{N}}
  \end{displaymath}
  
  \centerline{\animategraphics[loop,controls,width=\textwidth]{10}{exp/delaycircular}{0}{49}}

\end{frame}

\begin{frame}
  \frametitle{Delayed impulse is really periodic impulse}

  \begin{displaymath}
    \delta\left[(\!(n-n_0)\!)_N\right] \leftrightarrow e^{-j\frac{2\pi kn_0}{N}}
  \end{displaymath}
  
  \centerline{\animategraphics[loop,controls,width=\textwidth]{10}{exp/delayperiodic}{0}{49}}

\end{frame}

%%%%%%%%%%%%%%%%%%%%%%%%%%%%%%%%%%%%%%%%%%%%
\section[Phase]{Phase Unwrapping}
\setcounter{subsection}{1}


\begin{frame}
  \frametitle{Principal Phase}

  \begin{itemize}
  \item Something weird going on: how can the phase keep getting
    bigger and bigger, but the signal wraps around?
  \item It's because the phase wraps around too!
    \begin{displaymath}
      \measuredangle X[k] = -\omega_k (N+n) = -\omega_k n,~~~\mbox{because}~\omega_k=\frac{2\pi k}{N}
    \end{displaymath}
  \item {\bf Principal phase =} add $\pm 2\pi$ to the phase, as
    necessary, so that $-\pi< \measuredangle X[k]\le \pi$
  \item {\bf Unwrapped phase = } let the phase be as large as
    necessary so that it is plotted as a smooth function of $\omega$
  \end{itemize}
\end{frame}

\begin{frame}
  \frametitle{Unwrapped phase vs. Principal phase}

  \begin{displaymath}
    \delta\left[(\!(n-n_0)\!)_N\right] \leftrightarrow e^{-j\frac{2\pi kn_0}{N}}
  \end{displaymath}
  
  \centerline{\animategraphics[loop,controls,width=\textwidth]{10}{exp/principalphase}{0}{49}}

\end{frame}


%%%%%%%%%%%%%%%%%%%%%%%%%%%%%%%%%%%%%%%%%%%%
\section[MRI]{MRI}
\setcounter{subsection}{1}

\begin{frame}
  \frametitle{What is MRI?}

  To learn a little about spatial aliasing in MRI, let's start the first section of MP4.
\end{frame}

%%%%%%%%%%%%%%%%%%%%%%%%%%%%%%%%%%%%%%%%%%%%
\section[Summary]{Summary}
\setcounter{subsection}{1}

\begin{frame}
  \frametitle{Summary}

  Properties of the DFT:

  \begin{enumerate}
    \setcounter{enumi}{-1}
  \item Periodicity: $X[k+N]=X[k]$, $x[n+N]=x[n]$
  \item Conjugate symmetric: $x[n]$ Real $\leftrightarrow X[k]=X^*[-k]$
  \item Linearity:
    \[z[n]=ax[n]+by[n]\leftrightarrow Z[k]=aX[k]+bY[k]
    \]
  \item Time Shift: $x[n-n_0]\leftrightarrow e^{-j\frac{2\pi
      kn_0}{N}}X[k]$, paying attention to the fact that $x[n]$ is
    periodic in time
  \item Frequency Shift: $e^{j\frac{2\pi k_0 n}{N}}x[n]\leftrightarrow
    X[k-k_0]$, paying attention to the fact that $X[k]$ is periodic
    in frequency
  \end{enumerate}
\end{frame}

\end{document}
